% answers_short.tex —— 2008 年·普通高等学校招生考试·辽宁卷·理科 · 简短答案速查
% 由 tools/make_short_answers.py 自 answers.tex 的 \gksolhead/\gkanswer 提取，
% 并与 2008/2008_数学真题_辽宁理_答案要点.md 逐题比对后生成，勿手改（改源 MD 后重跑）。

\begin{gknotice}
\textbf{使用说明}：本节只给最终答案便于快速核对，完整推导见解析版。答案为编写者独立推导，\textbf{非官方标准答案}。
\end{gknotice}

\gkpart{选择题}
\begin{center}\begin{tabular}{ccccccccccccc}
\toprule
题号 & 1 & 2 & 3 & 4 & 5 & 6 & 7 & 8 & 9 & 10 & 11 & 12 \\
\midrule
答案 & D & B & C & B & A & A & C & A & B & A & D & C \\
\bottomrule
\end{tabular}\end{center}

\gkpart{填空题}
\begin{itemize}[leftmargin=2.6em]
  \item[\textbf{13.}] $y=\begin{cases}x-1,&x<1,\\ \ln x,&x\geqslant1.\end{cases}$
  \item[\textbf{14.}] $\dfrac32$。
  \item[\textbf{15.}] $5$。
  \item[\textbf{16.}] $\dfrac{14}3$。
\end{itemize}

\gkpart{解答题}
\begin{itemize}[leftmargin=2.6em]
  \item[\textbf{17.}] （1）$a=b=2$；（2）$\dfrac{2\sqrt3}3$。
  \item[\textbf{18.}] （1）$0.2$，$0.5$，$0.3$；（2）分布列如下，$E\xi=12.4$。
  \item[\textbf{19.}] （1）证明见解析；（2）定值 $\sqrt2$；（3）$\dfrac{\sqrt2}6$。
  \item[\textbf{20.}] （1）$x^2+\dfrac{y^2}4=1$；（2）$k=\pm\dfrac12$；（3）证明见解析。
  \item[\textbf{21.}] （1）$a_2=6$，$a_3=12$，$a_4=20$，$b_2=9$，$b_3=16$，$b_4=25$，$a_n=n(n+1)$，$b_n=(n+1)^2$；（2）证明见解析。
  \item[\textbf{22.}] （1）增区间 $(0,1)$，减区间 $(1,+\infty)$，极大值 $f(1)=\ln2$，无极小值；（2）不存在。
\end{itemize}
